\documentclass[10pt,a4paper]{amsart}
\usepackage[margin=19mm]{geometry}
\usepackage{amsmath,amssymb,mathtools}
\usepackage[scr=rsfs,cal=euler]{mathalfa}
\usepackage{xfrac}
\usepackage[unicode,hidelinks,bookmarks=false]{hyperref}
\newcommand{\ie}{i.e.\ }
\newcommand{\cf}{cf.\ }
\newcommand{\resp}{respectively\ }
\newcommand{\Subsection}{\subsection}
\providecommand{\nobreakdash}{\nobreak\hbox{-}}
\setlength{\emergencystretch}{2em}
\multlinegap0pt
\usepackage{bm}
\usepackage{bbm}
\usepackage{dsfont}
\usepackage{xspace}
\usepackage{cancel}
\usepackage{mathtools}
\usepackage{amsthm}
\makeatletter
\@ifundefined{theorem}{\newtheorem{theorem}{Theorem}}{}
\@ifundefined{lemma}{\newtheorem{lemma}[theorem]{Lemma}}{}
\@ifundefined{proposition}{\newtheorem{proposition}[theorem]{Proposition}}{}
\@ifundefined{remark}{\newtheorem{remark}[theorem]{Remark}}{}
\makeatother
\newcommand{\Findim}{\mathrm{Findim}\,}
\newcommand{\ddell}{\mathrm{ddell}\,}
\newcommand{\dell}{\mathrm{dell}\,}
\newcommand{\dirsum}{\xhookrightarrow{\!\oplus}}
\newcommand{\findim}{\mathrm{findim}\,}
\newcommand{\gldim}{\mathrm{gldim}\,}
\newcommand{\op}{\mathrm{op}}
\newcommand{\pd}{\mathrm{pd}\,}
\renewcommand{\k}{\mathbb{K}}
\renewcommand{\mod}{\mathrm{mod}\,}
\title[Homological Invariants of Left and Right Serial Quiver Algeb]{Homological Invariants of Left and Right Serial Quiver Algebras}
\author{Ruoyu Guo}
\date{Source: arXiv:2506.02307v2 (CC BY 4.0)}
\begin{document}
\maketitle

\noindent\emph{Source and licence.} Homological Invariants of Left and Right Serial Quiver Algebras. Version arXiv:2506.02307v2 (CC BY 4.0), distributed under the Creative Commons Attribution 4.0 International licence, as recorded in the arXiv licence metadata for this exact version and shown on the abstract page https://arxiv.org/abs/2506.02307v2. Full rights notice: licenses/arxiv-2506.02307-CC-BY-4.0.txt.\par\smallskip

\noindent\textit{Editorial scope.} Counted unit: the Theorem whose source label is \texttt{thm:right serial} in the article (source0.tex lines 438--444, source label \texttt{thm:right serial}), delivered together with the article's own complete original proof of it (source0.tex lines 446--479, one \texttt{proof} environment of 34 lines). No mathematics is written, repaired or re-derived: statement and proof are the article's own text line for line. Required context retained uncounted: eq:generic sequence of right completions (lines 348--351); eq:generic sequence of right completions factored (lines 353--356); eq:generic sequence of right completions reversed (lines 358--361); lem:existence of module with pd equal dell (lines 364--367); prop:tree dell findim (lines 382--388); eq:tree dell findim factored (lines 401--404); rmk:dell=2 case (lines 481--484). Every definition, display and label of those lines is retained unchanged. Omitted from this package and not redistributed: 1--347, 352--352, 357--357, 362--363, 368--381, 389--400, 405--437, 445--445, 480--480, 485--629. Nothing inside a retained excerpt is omitted or altered: every retained body is delivered exactly as the article's lines are written, so no omission receipt of any kind accompanies this package. Citations: the retained excerpts carry no citation command at all, so no bibliography entry of the article is transcribed and no reference is redistributed. Reference adaptation: none was needed and none was made; every reference inside the delivered unit resolves inside it, so no unresolved reference or citation is produced. Printed slips kept literal: no printed word, symbol, hypothesis or number of the counted statement or of its proof was altered. Layout and class: the article's own class is \texttt{conm-p-l} with packages including geometry, inputenc, babel, mathtools, graphicx, amsmath, amssymb, amsthm, comment, enumitem, leftidx, tikz, hyperref; the wrapper is the lane's pinned \texttt{amsart} editorial wrapper at 10pt on A4, which restates the theorem-like environments the retained text needs and adds the article's own macro definitions that the retained text needs (\texttt{\textbackslash Findim}, \texttt{\textbackslash ddell}, \texttt{\textbackslash dell}, \texttt{\textbackslash dirsum}, \texttt{\textbackslash findim}, \texttt{\textbackslash gldim}, \texttt{\textbackslash k}, \texttt{\textbackslash mod}, \texttt{\textbackslash op}, \texttt{\textbackslash pd}), with extra wrapper packages (bm, bbm, dsfont, xspace, cancel, mathtools). The delivered numbering is the unit's own. Assets: no retained span contains a figure environment, a graphics-inclusion command, a PDF-page inclusion, a table or a TikZ picture; no image file is copied into this package, the package declares no asset dependency and no image file is redistributed with it. Licence: version arXiv:2506.02307v2 of this article is distributed under the Creative Commons Attribution 4.0 International licence, as recorded in the arXiv licence metadata for this exact version and shown on the abstract page https://arxiv.org/abs/2506.02307v2. A full rights notice is delivered at licenses/arxiv-2506.02307-CC-BY-4.0.txt. Characters: every retained excerpt is pure ASCII, so the delivered engine, XeLaTeX with its default Latin Modern text font, reports no missing character.
\begin{equation}
\label{eq:generic sequence of right completions}
\tilde{\mu}=\widetilde{p_0}\widetilde{p_1}\cdots \widetilde{p_n},
\end{equation}

\begin{equation}
\label{eq:generic sequence of right completions factored}
\widetilde{x_0}\widetilde{y_0}\widetilde{x_1}\widetilde{y_1}\cdots \widetilde{x_{n-1}}\widetilde{y_{n-1}} \widetilde{p_n}.
\end{equation}

\begin{equation}
\label{eq:generic sequence of right completions reversed}
p_n y_{n-1}x_{n-1}\cdots y_1x_1y_0x_0.
\end{equation}

\begin{lemma}
\label{lem:existence of module with pd equal dell}
Using the notation in \eqref{eq:generic sequence of right completions}, \eqref{eq:generic sequence of right completions factored}, and \eqref{eq:generic sequence of right completions reversed}, there are at least $n$ minimal relations in \eqref{eq:generic sequence of right completions reversed}, and $\pd\!_{\Lambda} (\Lambda/p_n\Lambda) \geq n+1$.
\end{lemma}

\begin{proposition}
\label{prop:tree dell findim}
If $Q$ is an acyclic quiver and $\Lambda=\k Q/I$ is a monomial algebra, then 
\[
\gldim\Lambda=\Findim\Lambda=\findim\Lambda = \dell\Lambda^{\op}<\infty.
\]
\end{proposition}

\begin{equation}
\label{eq:tree dell findim factored}
\widetilde{y_0}\widetilde{x_1}\widetilde{y_1}\cdots \widetilde{x_{n-1}}\widetilde{y_{n-1}} \widetilde{p_n},
\end{equation}

\begin{remark}
\label{rmk:dell=2 case}
Note that in the proof of Theorem \ref{thm:right serial}, we did not need to explicitly prove the case $\dell\Lambda^{\op}=2$. In that case, the sequence of right completions in $Q^{\op}$ is simply $\widetilde{y_0}\widetilde{p_1}$. Using the same argument as in the proof, we can show $y_0\Lambda$ must be projective in $\mod\Lambda$ to ensure $\dell\Lambda^{\op} > 1$. Since the argument is essentially verbatim, we did not include it in the proof, but the same connection between the delooping level and projective dimension still exists for $\dell\Lambda^{\op}=2$.
\end{remark}

\begin{theorem}
\label{thm:right serial}
Let $\Lambda=\k Q/I$ be a right serial quiver algebra. Then the right little and big finitistic dimensions of $\Lambda$ are equal to the left delooping level and derived delooping level of $\Lambda$, \textit{i.e.},
\[
\findim\Lambda = \Findim\Lambda = \dell\Lambda^{\op} = \ddell\Lambda^{\op} < \infty.
\]
\end{theorem}

\begin{proof}
Since right serial quiver algebras are monomial, the four quantities are all finite. Also note that each vertex in $Q$ (resp. $Q^{\op}$) has at most one outgoing (resp. incoming) arrow. As in the proof of Proposition \ref{prop:tree dell findim}, it suffices to show $\dell\Lambda^{\op}\leq\findim\Lambda$. 

Let $S=S_{\tilde{v}}$ be a simple module in $\mod\Lambda^{\op}$ such that $\dell\!_{\Lambda^{\op}} S = \dell\Lambda^{\op}$. We know $\findim\Lambda=\Findim\Lambda=0$ if and only if $\dell\Lambda^{\op}=0$, so the theorem is true when $\dell\Lambda^{\op}$ is 0. The theorem is also true when $\dell\Lambda^{\op}=1$ since $\Findim\Lambda$ can only be 0 or 1 in that case, and $\Findim\Lambda=0$ implies $\dell\Lambda^{\op}=0$. We will prove the statement for $\dell\Lambda^{\op}\geq 3$ (corresponding to $n\geq 2$ in \eqref{eq:right serial right completion factored}), so the remaining $\dell\Lambda^{\op}=2$ case will follow, which we also explain in Remark \ref{rmk:dell=2 case}. Consider any arrow $\widetilde{y_0}$ starting at $\tilde{v}$ so that there exists a sequence of right completions
\begin{equation}
\label{eq:right serial right completion factored}
\widetilde{y_0}\widetilde{x_1}\widetilde{y_1}\widetilde{x_2}\widetilde{y_2} \cdots \widetilde{x_{n-1}}\widetilde{y_{n-1}}\widetilde{p_n}.
\end{equation}
The sequence \eqref{eq:right serial right completion factored} is factored in the same way as \eqref{eq:tree dell findim factored} in the proof of Proposition \ref{prop:tree dell findim}, where $\widetilde{y_{n-1}}\widetilde{p_n}$ and $\widetilde{y_i}\widetilde{x_{i+1}}\widetilde{y_{i+1}}$ for $i=0,1,\dots,n-2$ are minimal relations in $B_{\tilde{I}}$. Since $\dell\Lambda^{\op} = \dell_{\Lambda^{\op}} S = n+1 \geq 3$, we may choose \eqref{eq:right serial right completion factored} to correspond to the information that there exists a summand of the form $\widetilde{x_{n-1}}\widetilde{y_{n-1}}\Lambda^{\op}$ of $\Omega_{\Lambda^{\op}}^n S$ that is not $(n+1)$-deloopable and after we take another syzygy, the summand $\widetilde{p_n}\Lambda^{\op}$ of $\Omega_{\Lambda^{\op}}^{n+1} S$ is at least $(n+2)$-deloopable. We also know $\widetilde{y_0}\Lambda^{\op}\dirsum \Omega_{\Lambda^{\op}} S$ and $\widetilde{x_i}\widetilde{y_i}\Lambda^{\op}\dirsum \Omega_{\Lambda^{\op}}^{i+1} S$ for $1\leq i\leq n-1$.

Consider the reverse of \eqref{eq:right serial right completion factored} in $Q$ written as
\begin{equation}
\label{eq:right serial right completion reversed}
p_n y_{n-1}x_{n-1}\cdots y_1x_1y_0.
\end{equation}
We show that $x_1y_0\Lambda$ is always projective in $\mod\Lambda$. First, if $n=2$, $\dell\Lambda^{\op}=3$, but $x_1y_0\Lambda$ is not projective, then there exists a right completion $p$ for $x_1y_0$ in $Q$. We factor $x_1=z_1z_2$ such that $z_2y_0p$ is a minimal relation in $B_I$. Then we have a sequence of right completions $p_2y_1z_1z_2y_0p$ in $Q$, and its reverse in $Q^{\op}$ is
\begin{equation}
\label{eq:right serial dell=3 sequence}
\widetilde{p}\widetilde{y_0}\widetilde{z_2}\widetilde{z_1}\widetilde{y_1}\widetilde{p_2},
\end{equation}
where the minimal relations are $\widetilde{p}\widetilde{y_0}\widetilde{z_2}$, $\widetilde{y_0}\widetilde{z_2}\widetilde{z_1}\widetilde{y_1}$, and $\widetilde{y_1}\widetilde{p_2}$. So, we find that $\widetilde{z_1}\widetilde{y_1}\Lambda^{\op}$ is a summand of $\Omega^3_{\Lambda^{\op}} (\Lambda^{\op}/\widetilde{p}\Lambda^{op})$. On the other hand, $\widetilde{z_2}\widetilde{z_1}\widetilde{y_1}\Lambda^{\op}$ is a summand of $\Omega^2_{\Lambda^{\op}} S_{\tilde{v}}$. We showed in the proof of Lemma \ref{lem:existence of module with pd equal dell} that there is no minimal relation in \eqref{eq:right serial dell=3 sequence} starting from any arrow in $\widetilde{z_2}\widetilde{z_1}$, so the summands of $\widetilde{z_2}\widetilde{z_1}\widetilde{y_1}\Lambda^{\op}$ and $\widetilde{z_1}\widetilde{y_1}\Lambda^{\op}$ that are supported in the vertices in $\widetilde{p_2}$ are equal. We can apply this argument to any such sequence in the form \eqref{eq:right serial dell=3 sequence} where $\widetilde{y_0}$ is an arrow going out of $\tilde{v}$. This shows $\dell\!_{\Lambda^{\op}} S_{\tilde{v}} \leq 2$, contradicting the assumption.

Now we continue to the case $n\geq 3$ and $\dell\Lambda^{\op}\geq 4$. Suppose for a contradiction that $x_1y_0\Lambda$ is not projective. Then in the same way there exist a nonzero path $p$ and a factorization $x_1=z_1z_2$ such that $z_2y_0p$ is a minimal relation in $B_I$. Reversing to $Q^{\op}$, we get the sequence of right completions
\begin{equation}
\label{eq:right serial right completion factored additional term}
\widetilde{p}\widetilde{y_0}\widetilde{z_2}\widetilde{z_1}\widetilde{y_1}\widetilde{x_2}\widetilde{y_2} \cdots \widetilde{x_{n-1}}\widetilde{y_{n-1}}\widetilde{p_n}.
\end{equation}
Let $M$ be the non-projective summand of $\Lambda^{\op}/\tilde{p}\Lambda^{\op}\in\mod\Lambda^{\op}$. We get that the summand $\widetilde{x_2}\widetilde{y_2}\Lambda^{\op}$ of $\Omega_{\Lambda^{\op}}^3 S$ is a summand of the fourth syzygy of $M$. This also contradicts the assumption that $\dell\!_{\Lambda^{\op}} S\geq 4$. Therefore, $x_1y_0\Lambda$ is always projective in $\mod\Lambda$.

Let $N\in\mod\Lambda$ be the non-projective summand of $\Lambda/p_n\Lambda$. Note that since $\Lambda$ is right serial, there is no other sequence of right completions starting with any arrow in \eqref{eq:right serial right completion reversed} except for the sequence \eqref{eq:right serial right completion reversed} itself. From the relations in \eqref{eq:right serial right completion reversed}, we find that $\Omega_{\Lambda}^n N = x_2y_1\Lambda$ is not projective, but $\Omega_{\Lambda}^{n+1} N= x_1y_0\Lambda$ is projective. Thus, we complete the proof because
\[
\findim\Lambda \geq \pd\!_{\Lambda} N = n+1 = \dell\Lambda^{\op}\geq\ddell\Lambda^{\op}\geq\Findim\Lambda\geq\findim\Lambda.
\]
\end{proof}

\end{document}
